% =============================================================================
% Scrivia — Argument writing
% Modello di progetto LaTeX (libro o saggio argomentativo)
% =============================================================================
%
% Struttura dei file:
%   main.tex            questo file: metadati, stile e ordine dei capitoli
%   misc/options.tex    stile tipografico (font, pagina, titoli, riquadri, indici)
%   frontmatter/        copertina, prefazione, retrocopertina
%   content/            i capitoli: un file .tex per capitolo
%   images/             le immagini (\includegraphics cerca qui)
%
% L'ordine dei capitoli è quello delle righe \input{...} qui sotto: lo Studio
% LaTeX di Scrivia lo legge per costruire la struttura del progetto. Il commento
% dopo \input diventa il titolo mostrato nell'elenco.
% =============================================================================

% --- Modalità bozza / finale -------------------------------------------------
% \isdrafttrue  = BOZZA (filigrana, immagini come segnaposto)
% \isdraftfalse = FINALE
% Il workflow "Genera PDF" cambia questa riga da solo in base alla scelta.
\newif\ifisdraft
\isdraftfalse   % <<< CAMBIA QUI: true = BOZZA, false = FINALE

\ifisdraft
  \PassOptionsToPackage{draft}{graphicx}
\fi

\documentclass[11pt,twoside]{book}

\usepackage[T1]{fontenc}
\usepackage[utf8]{inputenc}
\usepackage{textcomp}
\usepackage{xcolor}
\usepackage[outline]{contour}
\contourlength{1.1pt}
\usepackage{comment}

% --- Metadati -----------------------------------------------------------------
\newcommand{\authorname}{Nome Cognome}
\newcommand{\booktitle}{Titolo del saggio}
\newcommand{\subtitleA}{Una tesi chiara,}
\newcommand{\subtitleB}{sostenuta da argomenti e prove}
\newcommand{\booksubject}{Saggio argomentativo}
\newcommand{\bookkeywords}{argomentazione, tesi, prove}

\newcommand{\publisher}{Editore}
\newcommand{\editionyear}{\the\year}
\newcommand{\isbn}{978-00-00-00000-0}

% --- Versione della bozza -------------------------------------------------------
\newcommand{\draftversion}{0.1}
\newcommand{\draftdate}{\today}

% Stile e opzioni del modello
% =============================================================================
% Modern Professional Book Style Package
% Updated with Source Serif Pro and contemporary design
% Modello di Scrivia — Argument writing (stile derivato da libro-futuro)
% =============================================================================

% Font encoding first
\usepackage[T1]{fontenc}
\usepackage[utf8]{inputenc}

% Microtype for superior typography
\usepackage[
activate={true,nocompatibility},
final,
tracking=true,
kerning=true,
spacing=true,
factor=1100,
stretch=10,
shrink=10
]{microtype}

% LANGUAGE SETTINGS ------------------------------------------------------
\usepackage[english,italian]{babel}
\addto\captionsitalian{%
	\renewcommand{\contentsname}{Indice}
	\renewcommand{\chaptername}{Capitolo}
}

% MODERN FONT SETUP - Source Serif Pro -----------------------------------
\usepackage[default,semibold]{sourceserifpro}
\usepackage{sourcesanspro}  % For headers and display text
\usepackage{sourcecodepro}  % For any code snippets

% Optimal line spacing for Source Serif
\linespread{1.18}

% TYPOGRAPHY ENHANCEMENTS ------------------------------------------------
\usepackage{lettrine}   % Drop caps
\usepackage{graphicx}   
\usepackage{xcolor}
\usepackage{tcolorbox} % Per creare riquadri personalizzabili
\tcbuselibrary{skins, breakable} % Librerie aggiuntive per i riquadri

% Modern color palette
\definecolor{modernblack}{RGB}{25, 25, 25}      % Softer than pure black
\definecolor{modernred}{RGB}{215, 38, 61}       % Contemporary accent
\definecolor{moderngray}{RGB}{110, 110, 110}    % For secondary elements
\definecolor{lightgray}{RGB}{245, 245, 245}     % For backgrounds
\definecolor{darkgray}{RGB}{70, 70, 70}         % For strong emphasis

% Professional paragraph settings
\setlength{\parskip}{0pt}                   
\setlength{\parindent}{1.2em}              
\frenchspacing
\widowpenalty=10000
\clubpenalty=10000
\raggedbottom  % Prevents ugly vertical stretching

% PAGE LAYOUT (17x24 cm - Saggio Grande) ---------------------------------
\usepackage[
paperwidth=170mm,
paperheight=240mm,
top=20mm,
bottom=25mm,
outer=20mm,
inner=25mm,
headsep=8mm,
footskip=12mm
]{geometry}

% Better line breaking
\emergencystretch=3em
\hfuzz=.5pt  % Suppress trivial overfull warnings

% TABLE OF CONTENTS STYLING ----------------------------------------------
\usepackage[titles]{tocloft}
\renewcommand{\cftchapfont}{\sffamily\normalsize}
\renewcommand{\cftchappagefont}{\sffamily\normalsize\color{moderngray}}
\renewcommand{\cftchapleader}{\cftdotfill{\cftdotsep}}
\renewcommand{\cftdotsep}{2.0}
\setlength{\cftbeforechapskip}{0.5em}

% CHAPTER STYLING --------------------------------------------------------
\usepackage{titlesec}

% Modern chapter style with sans-serif
\titleformat{\chapter}[display]
{\normalfont\sffamily\centering}
{\huge\textls[50]{\MakeUppercase{\chaptertitlename}}\quad\thechapter}
{25pt}
{\Huge\bfseries}

\titlespacing*{\chapter}{0pt}{40pt}{60pt}

% Section and subsection styles
\titleformat{\section}
{\normalfont\sffamily\Large\bfseries\color{darkgray}}
{\thesection}
{1em}
{}

\titleformat{\subsection}
{\normalfont\sffamily\large\bfseries\color{darkgray}}
{\thesubsection}
{1em}
{}

\titleformat{\subsubsection}
{\normalfont\sffamily\normalsize\bfseries\color{moderngray}}
{\thesubsubsection}
{1em}
{}

% HEADERS AND FOOTERS ---------------------------------------------------
\usepackage{fancyhdr}
\setlength{\headheight}{25pt}

% Main matter page style
\fancypagestyle{mainmatter}{%
	\fancyhf{}
	% Even pages: page number and author
	\fancyhead[LE]{\sffamily\footnotesize\color{moderngray}\thepage\quad\textls[50]{\MakeUppercase{\authorname}}}
	% Odd pages: chapter title and page number
	\fancyhead[RO]{\sffamily\footnotesize\color{moderngray}\textls[50]{\MakeUppercase{\leftmark}}\quad\thepage}
	% No header rule for cleaner look
	\renewcommand{\headrulewidth}{0pt}
}

% Plain style for chapter openings
\fancypagestyle{plain}{%
	\fancyhf{}
	\fancyfoot[C]{\sffamily\small\color{moderngray}\thepage}
	\renewcommand{\headrulewidth}{0pt}
}

% QUOTATION ENVIRONMENT --------------------------------------------------
\usepackage{csquotes}
\renewenvironment{quotation}
{\list{}{\listparindent 0em%
		\itemindent    0em
		\leftmargin    2em
		\rightmargin   2em
		\parsep        0.5em
		\topsep        0.5em}%
	\item\relax\small\color{darkgray}}
{\endlist}

% Block quotes with modern style
\newenvironment{modernquote}
{\begin{quotation}\sffamily\itshape}
	{\end{quotation}}

% LISTS - Modern Spacing -------------------------------------------------
\usepackage{enumitem}
\setlist[itemize]{
	topsep=0.5em,
	itemsep=0.25em,
	parsep=0.25em,
	leftmargin=1.5em,
	label={\color{modernred}\textbullet}
}
\setlist[enumerate]{
	topsep=0.5em,
	itemsep=0.25em,
	parsep=0.25em,
	leftmargin=1.5em
}

% CAPTIONS ---------------------------------------------------------------
\usepackage[
font={scriptsize,sf},
labelfont={bf,color=darkgray},
textfont={color=moderngray},
justification=centering,
margin=0.5cm
]{caption}
\setlength{\abovecaptionskip}{10pt}
\setlength{\belowcaptionskip}{10pt}

% FOOTNOTES --------------------------------------------------------------
\usepackage[bottom,hang,flushmargin]{footmisc}
\setlength{\footnotesep}{12pt}
\renewcommand{\footnotelayout}{\footnotesize\color{moderngray}}

% CUSTOM COMMANDS --------------------------------------------------------
% First few words in sans-serif bold
\newcommand{\firstwords}[1]{{\sffamily\bfseries #1}}

% Small caps with tracking
\newcommand{\smallcaps}[1]{\textsc{\textls[50]{#1}}}

% Book title formatted
\newcommand{\booktitleformatted}{{\sffamily\textit{\booktitle}}}

% Chapter quotes
\newcommand{\chapterquote}[2]{%
	\vspace{-0.5em}
	\begin{flushright}
		\begin{minipage}{0.8\textwidth}
			\raggedleft
			\sffamily\small\itshape\color{moderngray}
			``#1''\\[0.5em]
			\sffamily\footnotesize\color{darkgray}#2
		\end{minipage}
	\end{flushright}
	\vspace{1em}
}

% Modern lettrine settings
\renewcommand{\LettrineFontHook}{\sffamily\bfseries\color{modernblack}}
\setlength{\DefaultNindent}{0pt}
\setlength{\DefaultSlope}{0pt}
\renewcommand{\LettrineTextFont}{\sffamily}
\setlength{\DefaultFindent}{0.15em}

% PAGE NUMBERS WITH STYLE ------------------------------------------------
\usepackage{lastpage}

% HYPERREF (load last) ---------------------------------------------------
\usepackage[
colorlinks=true,
linkcolor=modernblack,
citecolor=modernred,
urlcolor=modernred,
bookmarks=true,
bookmarksopen=true,
bookmarksnumbered=true,
pdfauthor={\authorname},
pdftitle={\booktitle},
pdfsubject={\booksubject},
pdfkeywords={\bookkeywords},
pdfproducer={LaTeX with Source Serif Pro},
pdfcreator={pdfLaTeX}
]{hyperref}

% URL STYLE
\urlstyle{same}  % URLs in same font as text

% ADDITIONAL MODERN TOUCHES ----------------------------------------------
% Better spacing around display math (if used)
\usepackage{amsmath}
\AtBeginDocument{%
	\setlength\abovedisplayskip{12pt plus 3pt minus 3pt}%
	\setlength\belowdisplayskip{12pt plus 3pt minus 3pt}%
}

% Epigraphs for chapter quotes
\usepackage{epigraph}
\setlength{\epigraphwidth}{0.7\textwidth}
\renewcommand{\epigraphflush}{flushleft}
\renewcommand{\sourceflush}{flushleft}

% Optional: Add TikZ for any diagrams
\usepackage{tikz}
\usetikzlibrary{shapes,arrows,positioning}

% Optional: Better tables
\usepackage{booktabs}
\renewcommand{\arraystretch}{1.2}
% =============================================================================
% INDEXING PACKAGE AND CONFIGURATION
% =============================================================================
\usepackage[xindy]{imakeidx}

% Stile professionale per l'indice
\indexsetup{
	othercode=\small,
	firstpagestyle=plain,
	headers={Indice analitico}{Indice analitico}
}

% Indice analitico principale
\makeindex[
name=main,
title=Indice analitico,
program=texindy,
options=-L italian -C utf8 -I latex,
columns=2 
]

% Secondo indice (facoltativo): concetti chiave
\makeindex[
name=concetti,
title=Indice dei concetti,
program=texindy,
options=-L italian -C utf8 -I latex,
columns=2 
]

% =============================================================================
% STILI PER RIQUADRI CONCETTUALI E CODICE
% =============================================================================

% --- Riquadro per concetti chiave ---
\newtcolorbox{concetto}[1][]{
	colback=lightgray,       
	colframe=darkgray,       
	fonttitle=\sffamily\bfseries,
	coltitle=modernblack,
	breakable,
	enhanced,
	arc=3mm,
	boxrule=0.8pt,
	title=#1
}

% --- Riquadro per lo pseudo-codice ---
\newtcolorbox{codice}[1][]{
	colback=darkgray,       
	coltext=white,
	fontupper=\ttfamily,      
	breakable,
	arc=1mm,
	boxrule=0pt,
	title=#1,
	fonttitle=\sffamily\bfseries,
	coltitle=white
}

% --- Riquadro per Dati e Grafici ---
\newtcolorbox{boxdati}[1][]{
	colback=white,            
	colframe=modernred,       
	fonttitle=\sffamily\bfseries,
	coltitle=white,
	breakable,
	enhanced,
	arc=3mm,
	boxrule=0.8pt,
	title=#1
}

% =============================================================================
% STILI PER DIALOGO UMANO-AI
% =============================================================================

\definecolor{humancolor}{RGB}{25, 25, 25}      
\definecolor{aicolor}{RGB}{70, 130, 180}       

\newcommand{\humanbox}[1]{%
	\vspace{0.8em}%
	\begin{tcolorbox}[
		colback=white,
		colframe=humancolor,
		leftrule=3pt,
		rightrule=0pt,
		toprule=0pt,
		bottomrule=0pt,
		left=10pt,
		right=5pt,
		top=5pt,
		bottom=5pt
		]
		\noindent\textbf{\sffamily\color{humancolor}UMANO:} #1
	\end{tcolorbox}%
	\vspace{0.3em}%
}

\newcommand{\aibox}[1]{%
	\vspace{0.8em}%
	\begin{tcolorbox}[
		colback=aicolor!5,
		colframe=aicolor,
		leftrule=3pt,
		rightrule=0pt,
		toprule=0pt,
		bottomrule=0pt,
		left=10pt,
		right=5pt,
		top=5pt,
		bottom=5pt
		]
		\noindent\textbf{\sffamily\color{aicolor}AI:} #1
	\end{tcolorbox}%
	\vspace{0.3em}%
}

% Immagine con segnaposto: se il file non esiste ancora in images/ compare un
% riquadro con la descrizione, e il libro compila lo stesso. Quando il file
% viene aggiunto, l'immagine appare da sola.
% Uso: \immagine{images/nome.png}{descrizione di cosa deve mostrare}
% Dimensionamento delle immagini: tutta la larghezza del testo, ma mai piu'
% alte di 0.87 dell'altezza utile, cosi' immagine e didascalia stanno sempre
% nella stessa pagina. Un'immagine alta finisce da sola in una pagina che
% riempie quasi per intero; una bassa si incastra nel testo senza vuoti.
\setkeys{Gin}{keepaspectratio,height=0.87\textheight}
\newcommand{\immagine}[2]{%
	\IfFileExists{#1}%
		{\includegraphics[width=\textwidth]{#1}}%
		{\fbox{\parbox{0.8\textwidth}{\centering\small\vspace{1.5em}%
			[IMMAGINE DA INSERIRE: #2]\\[0.4em]\texttt{\detokenize{#1}}\vspace{1.5em}}}}%
}

\usepackage{graphicx}
\usepackage{eso-pic}
\usepackage{fancyhdr}
\graphicspath{{images/}}

\usepackage{ragged2e}
\justifying

% --- Intestazioni: diverse per BOZZA e FINALE --------------------------------
\ifisdraft
  \usepackage{draftwatermark}
  \SetWatermarkText{BOZZA v\draftversion}
  \SetWatermarkScale{0.5}
  \SetWatermarkColor[gray]{0.95}
  \SetWatermarkAngle{45}
  \fancypagestyle{mainmatter}{%
    \fancyhf{}
    \fancyhead[LE]{\sffamily\footnotesize\color{moderngray}\thepage\quad\textls[50]{\MakeUppercase{\chaptertitlename\ \thechapter}}}
    \fancyhead[RO]{\sffamily\footnotesize\color{moderngray}\textls[50]{\MakeUppercase{\leftmark}}\quad\thepage\quad\sffamily\tiny\color{modernred}[BOZZA~v\draftversion, \draftdate]}
    \renewcommand{\headrulewidth}{0pt}
  }
\else
  \fancypagestyle{mainmatter}{%
    \fancyhf{}
    \fancyhead[LE]{\sffamily\footnotesize\color{moderngray}\thepage\quad\textls[50]{\MakeUppercase{\chaptertitlename\ \thechapter}}}
    \fancyhead[RO]{\sffamily\footnotesize\color{moderngray}\textls[50]{\MakeUppercase{\leftmark}}\quad\thepage}
    \renewcommand{\headrulewidth}{0pt}
  }
\fi
\usepackage{bookmark}

\begin{document}

% ===== FRONT MATTER =====
\frontmatter
\pagestyle{empty}

% =============================================================================
% Copertina: titolo, sottotitolo e autore su un riquadro semitrasparente.
% Se esiste images/Cover.png viene usata come sfondo a tutta pagina.
% =============================================================================
\begin{titlepage}
	\thispagestyle{empty}
	\IfFileExists{images/Cover.png}{%
		\AddToShipoutPictureBG*{%
			\AtPageLowerLeft{\includegraphics[width=\paperwidth,height=\paperheight]{Cover.png}}%
		}%
	}{}
	\begin{tikzpicture}[remember picture,overlay]
		\node[
		anchor=north,
		align=center,
		yshift=-2.0cm,
		text=white,
		fill=modernblack,
		fill opacity=.85,
		text opacity=1,
		rounded corners=10pt,
		inner xsep=18pt,
		inner ysep=14pt,
		text width=.82\paperwidth
		] at (current page.north) {%
			{\fontsize{30}{36}\selectfont\sffamily\bfseries\textls[100]{\MakeUppercase{\booktitle}}\par}
			\vspace{0.8em}
			{\normalsize\sffamily\itshape \subtitleA\\[0.6em]\subtitleB\par}
			\vspace{0.9em}
			{\large\sffamily\textls[80]{\MakeUppercase{\authorname}}}
		};
		\node[
		anchor=south,
		fill=modernblack,
		text=white,
		rounded corners=6pt,
		inner xsep=10pt,
		inner ysep=6pt
		] at ([yshift=6mm]current page.south) {\large\sffamily\textls[50]{\MakeUppercase{\publisher}}};
	\end{tikzpicture}
\end{titlepage}
       % Copertina
\cleardoublepage
% =============================================================================
% Pagina del copyright
% =============================================================================
\thispagestyle{empty}
\vspace*{\fill}

\noindent
Copyright \textcopyright\ \editionyear\ \authorname

\vspace{0.5em}

\noindent
\small
Tutti i diritti riservati. Nessuna parte di questa pubblicazione può essere
riprodotta, memorizzata in sistemi di recupero o trasmessa in qualsiasi forma
o con qualsiasi mezzo senza il previo consenso scritto dell'editore.

\vspace{1em}

\noindent
Prima edizione: \editionyear\\
ISBN: \isbn\\
Pubblicato da: \publisher

\normalsize
   % Copyright
\cleardoublepage

\pagestyle{plain}
\tableofcontents
\cleardoublepage

% =============================================================================
% Prefazione
% =============================================================================
\chapter*{Prefazione}
\addcontentsline{toc}{chapter}{Prefazione}

Perché questo testo esiste, a chi si rivolge e che cosa il lettore troverà
alla fine del percorso. Una prefazione argomentativa anticipa la domanda a cui
il libro risponde, senza ancora dimostrare la risposta.
   % Prefazione

% ===== MAIN CONTENT =====
\mainmatter
\pagestyle{mainmatter}

% ========================================
% PARTE I: LA TESI
% ========================================
\part{La tesi}
% =============================================================================
% Introduzione: il problema e la tesi
% =============================================================================
\chapter{Introduzione}

\chapterquote{Una citazione che apre il problema.}{Autore, \textit{Opera}}

\lettrine{O}{gni testo argomentativo} parte da un problema concreto e prende
posizione. Qui si presenta il contesto, si spiega perché la questione conta e
si enuncia la \textbf{tesi} in modo chiaro e verificabile.\index{tesi}

\section{Il problema}
Descrivi la situazione e la domanda a cui il testo risponde.\footnote{Le note
a piè di pagina servono per fonti e precisazioni che non interrompono il
ragionamento.}

\section{La tesi}
\begin{concetto}[Tesi]
Una frase che il resto del testo dovrà dimostrare.
\end{concetto}

\section{Il percorso}
\begin{itemize}
	\item Primo argomento a sostegno.
	\item Secondo argomento a sostegno.
	\item Obiezioni principali e risposta.
\end{itemize}
   % Introduzione: il problema e la tesi

% ========================================
% PARTE II: GLI ARGOMENTI
% ========================================
\part{Gli argomenti}
% =============================================================================
% Gli argomenti a sostegno della tesi
% =============================================================================
\chapter{Gli argomenti}

\section{Primo argomento}
Affermazione, ragionamento e prova: un dato, un esempio, una fonte
autorevole.\index{prova}

\begin{quote}
Una citazione che sostiene l'argomento, con la sua fonte.
\end{quote}

\section{Secondo argomento}
Ogni argomento va collegato esplicitamente alla tesi: \emph{perché} questa
prova la rende più credibile?

% Per inserire un'immagine (il file va in images/):
% \begin{figure}[htbp]
%     \centering
%     \includegraphics[width=0.8\textwidth]{images/grafico.png}
%     \caption{Didascalia della figura}
%     \label{fig:grafico}
% \end{figure}
      % Gli argomenti a sostegno
% =============================================================================
% Obiezioni e confutazione
% =============================================================================
\chapter{Obiezioni e confutazione}

\section{L'obiezione più forte}
Riporta l'obiezione nella sua forma migliore, senza caricature.

\section{La risposta}
Mostra perché l'obiezione non regge, o in che misura la tesi va precisata.

\begin{boxdati}[Dati a confronto]
Numeri, fonti e limiti delle prove usate nella confutazione.
\end{boxdati}
      % Obiezioni e confutazione

% ========================================
% PARTE III: LA CONCLUSIONE
% ========================================
\part{La conclusione}
% =============================================================================
% Conclusione
% =============================================================================
\chapter{Conclusione}

Riprendi la tesi alla luce degli argomenti e delle obiezioni discusse, indica
le conseguenze pratiche e le domande che restano aperte.
    % Conclusione

% ===== BACK MATTER =====
\backmatter
\pagestyle{empty}

\cleardoublepage
\phantomsection
\addcontentsline{toc}{chapter}{Indice analitico}
\printindex[main]

% Secondo indice (facoltativo): togli il % se usi \index[concetti]{...}
% \cleardoublepage
% \phantomsection
% \addcontentsline{toc}{chapter}{Indice dei concetti}
% \printindex[concetti]

\cleardoublepage
% =============================================================================
% Retrocopertina
% =============================================================================
\pagestyle{empty}
\cleardoublepage
\begin{tikzpicture}[remember picture,overlay]
	\fill[modernblack] (current page.south west) rectangle (current page.north east);
	\node[anchor=north, text width=0.8\textwidth, align=center]
	at ([yshift=-2cm] current page.north) {%
		\color{white}
		{\sffamily\Large\textit{``La tesi del libro, in una frase.''}}\\[0.8em]
		{\sffamily\normalsize\textls[50]{\MakeUppercase{\authorname}}}
	};
	\node[anchor=center, text width=0.8\textwidth, align=justify]
	at (current page.center) {%
		\sffamily\small\color{white}
		Il problema, la tesi e i tre argomenti principali in poche righe:
		la quarta di copertina è un piccolo testo argomentativo.
	};
\end{tikzpicture}
        % Retrocopertina

\end{document}
